%% 3-0 序章　はかったデータから、きまりを見つけよう（第12回：組 V3-1）  % noruby
%% はかる → 表 → 点 → 線 → 式 の流れ、点は線にぴったりのらない（3-16 への宿題）、路線と駅、できるようになること、読み方。  % noruby
%% 2-15 の F=kx は教えなおさない（思い出すだけ）。  % noruby
\makeatletter
\gdef\bk@partlabel{\R{第}{だい}3\R{巻}{かん}}\gdef\bk@partname{\R{関数}{かん|すう}とデータ}
%% ---- この章だけで使う絵の部品（頭文字 \bkthra…） ----
% ばね（たて）  \bkthraspring{(x,上のはし)}{全体の長さ}  % noruby
\newcommand\bkthraspring[2]{\begin{scope}[shift={#1}]
  \draw[bkGray!90!black, line width=0.7pt] (0,0) -- (0,-2);
  \draw[bkGray!90!black, line width=0.7pt, decorate, decoration={coil, aspect=0.45, segment length=1.3mm, amplitude=2mm}] (0,-2) -- (0,{-(#2)+2});
  \draw[bkGray!90!black, line width=0.7pt] (0,{-(#2)+2}) -- (0,{-(#2)});
\end{scope}}
% おもり  \bkthraweight{(x,上のはし)}  はば7mm・高さ6mm  % noruby
\newcommand\bkthraweight[1]{\begin{scope}[shift={#1}]
  \filldraw[fill=bkGray!45, draw=bkGray!90!black, line width=0.6pt, rounded corners=0.8mm] (-3.5,-6) rectangle (3.5,0);
  \fill[white, opacity=0.35] (-2.6,-1.6) rectangle (-1.4,-4.8);
\end{scope}}
% 天井の板  \bkthraceil{(x1,y)}{(x2,y)}  % noruby
\newcommand\bkthraceil[2]{\fill[bkGray!35] #1 rectangle ($#2+(0,2.2)$);
  \draw[bkGray!90!black, line width=0.8pt] #1 -- #2;}
% 路線の1行  \bkthraroute{y}{路線番号}{前の巻}{この巻の駅}{この先}  % noruby
\newcommand\bkthraroute[5]{%
  \draw[bkRoute#2, line width=3pt, line cap=round] (10,#1) -- (132,#1);
  \node[inner sep=0pt] at (4,#1) {\bkrouteicon[7mm]{#2}};
  \fill[white] (24,#1) circle (2.2);\draw[bkRoute#2, line width=1.1pt] (24,#1) circle (2.2);
  \node[anchor=north, font=\sffamily\footnotesize, text=bkGray, inner sep=0pt] at (24,#1-3.8) {#3};
  \fill[white, rounded corners=1.5mm] (38,#1-4.4) rectangle (110,#1+4.4);
  \draw[bkRoute#2, line width=1.6pt, rounded corners=1.5mm] (38,#1-4.4) rectangle (110,#1+4.4);
  \node[font=\sffamily\footnotesize, text=bkInk] at (74,#1) {#4};
  \fill[white] (123,#1) circle (2.2);\draw[bkRoute#2, line width=1.1pt] (123,#1) circle (2.2);
  \node[anchor=north, font=\sffamily\footnotesize, text=bkGray, inner sep=0pt] at (123,#1-3.8) {#5};}
% 読み方のカード  \bkthracard{(x,y)}{アイコン}{色}{見出し}{中身}  % noruby
\newcommand\bkthracard[5]{%
  \node[anchor=north west, draw=#3, line width=1pt, fill=white, rounded corners=2.5mm,
    minimum width=64mm, minimum height=60mm, inner sep=0pt] at #1 {};
  \node[anchor=north west, inner sep=0pt] (ci) at ($#1+(3.5,-3)$) {\bkicon[1.5em]{#2}{#3}};
  \node[anchor=west, font=\sffamily\bfseries\small, text=#3!85!black, inner sep=0pt] at ([xshift=2mm]ci.east) {#4};
  \node[anchor=north west, text width=57mm, align=left, font=\small, inner sep=0pt] at ($#1+(3.5,-12)$) {#5};}
% 箇条書きの1項目（・で始め、2行目は字下げ）  \bkthraitem{文}  % noruby
\newcommand\bkthraitem[1]{\par\noindent\hangindent=1\zw\hangafter=1 ・#1\par}
% 章番号の札  \bkthrapill{章}  % noruby
\newcommand\bkthrapill[1]{\bk@pill{#1}}
% 会話の1行  \bkthratalk{y}{sora|mina|delta}{名前}{せりふ}  % noruby
\newcommand\bkthratalk[4]{%
  \bkChar{#2}{(7,#1-6.5)}{12mm}
  \node[rectangle callout, callout absolute pointer={(15.5,#1)}, anchor=west, draw=bkTheme, line width=0.8pt,
    fill=white, rounded corners=2.5mm, inner xsep=3mm, inner ysep=1.8mm, text width=103mm, align=left,
    font=\small] at (21,#1) {#4};
  \node[anchor=north, font=\sffamily\scriptsize, text=bkGray, inner sep=0pt] at (7,#1-7.5) {#3};}
% 流れの図の1まい  \bkthrastep{名前}{中心x}{見出し}  はば21mm・高さ30mm（y は -15〜15）  % noruby
\newcommand\bkthrastep[3]{%
  \node[draw=bkTheme, line width=1pt, fill=white, rounded corners=2mm, minimum width=21mm, minimum height=30mm, inner sep=0pt] (#1) at (#2,0) {};
  \node[anchor=north, font=\sffamily\bfseries\footnotesize, text=bkDeep, inner sep=0pt] at (#2,13) {#3};}
% 見出しのアイコン（路線：線と2つの駅。黄の地なので線は濃い紺）  \bksection{bkthraroute}{…} で使う  % noruby
\newcommand\bk@icon@bkthraroute{%
  \draw[bkOn, line width=0.13*\bk@iu, line cap=round] (-0.62,0) -- (0.62,0);
  \filldraw[fill=white, draw=bkOn, line width=0.09*\bk@iu] (-0.36,0) circle (0.2) (0.36,0) circle (0.2);}
\makeatother

\bkchapter{3}{0}{はかったデータから、\\きまりを\R{見}{み}つけよう}{\bklevel{2}}{はかったデータから、きまりを\R{見}{み}つけよう}

\noindent ばねにおもりをつるすと、ばねはのびます。おもりをふやすと、もっとのびます。ここまでは、だれでも\R{知}{し}っています。

では、おもりを1\R{個}{こ}ふやすごとに、\R{何}{なん}cm のびるのでしょう。10\R{個}{こ}つるしたら、\R{何}{なん}cm のびるのでしょう。「たくさんつるすと、たくさんのびる」ということばだけでは、\R{答}{こた}えられません。

\begin{hatena}
\noindent\begin{minipage}[c]{0.6\linewidth}
ばねにおもりを1\R{個}{こ}ずつ\R{下}{さ}げていくと、のびはどうなるかな？\par
はかった\R{数}{かず}を\R{表}{ひょう}とグラフにしたら、のびを\R{計算}{けい|さん}する\R{式}{しき}を\R{作}{つく}れるかな？
\end{minipage}\hfill
\begin{minipage}[c]{0.36\linewidth}\centering
\begin{tikzpicture}[x=1mm, y=1mm]
  \bkthraceil{(-1,0)}{(35,0)}
  \bkthraspring{(3,0)}{14}
  \bkthraspring{(15,0)}{17}\bkthraweight{(15,-17)}
  \bkthraspring{(28,0)}{20}\bkthraweight{(28,-20)}\bkthraweight{(28,-26.5)}
  \node[bk tick, anchor=north] at (3,-35) {0\R{個}{こ}};
  \node[bk tick, anchor=north] at (15,-35) {1\R{個}{こ}};
  \node[bk tick, anchor=north] at (28,-35) {2\R{個}{こ}};
  \node[font=\sffamily\bfseries\Large, text=bkOrange!80!black] at (36.5,-16) {?};
\end{tikzpicture}
\end{minipage}
\end{hatena}

この\R{巻}{かん}は、このような「はかった\R{数}{かず}から、きまりを\R{見}{み}つける」\R{方法}{ほう|ほう}を\R{学}{まな}ぶ\R{巻}{かん}です。\R{序章}{じょ|しょう}では、ばねの\R{実験}{じっ|けん}を1つ\R{最後}{さい|ご}まで\R{通}{とお}してやってみて、この\R{巻}{かん}の\R{旅}{たび}の\R{地図}{ち|ず}を\R{作}{つく}ります。

\begin{omoidasou}
\bkrecall{2-3}{\R{変}{か}わった\R{量}{りょう}は「あと－まえ」、1あたりの\R{量}{りょう}はわり\R{算}{ざん}で\R{求}{もと}める。}
\bkrecall{2-6}{はかった\R{値}{あたい}にはずれがある。1\R{目}{め}もりの10\R{分}{ぶん}の1まで\R{読}{よ}む。}
\bkrecall{2-15}{\R{比例}{ひ|れい} $y=ax$ のグラフは、\R{原点}{げん|てん}を\R{通}{とお}る\R{直線}{ちょく|せん}。ばねは $F=kx$。}
\end{omoidasou}

%% ---------------- はかる → 表 → 点 → 線 → 式 ----------------
\bksection{doguu}{はかる → \R{表}{ひょう} → \R{点}{てん} → \R{線}{せん} → \R{式}{しき}}

\R{科学者}{か|がく|しゃ}や\R{技術者}{ぎ|じゅつ|しゃ}は、\R{新}{あたら}しいきまりを\R{見}{み}つけるとき、いつも\R{同}{おな}じ5つの\R{段}{だん}を\R{進}{すす}みます。この\R{巻}{かん}の\R{章}{しょう}は、どれもこの5つの\R{段}{だん}のどこかをくわしくする\R{章}{しょう}です。

\begin{bkfigure}
\begin{tikzpicture}[x=1mm, y=1mm]
  \bkthrastep{s1}{10.5}{①\,はかる}
  \bkthrastep{s2}{38.25}{②\,\R{表}{ひょう}}
  \bkthrastep{s3}{66}{③\,\R{点}{てん}}
  \bkthrastep{s4}{93.75}{④\,\R{線}{せん}}
  \bkthrastep{s5}{121.5}{⑤\,\R{式}{しき}}
  \foreach \a/\b in {s1/s2,s2/s3,s3/s4,s4/s5} \draw[bk link] (\a) -- (\b);
  % ① ばねと目もり板
  \fill[bkGray!35] (4,5.5) rectangle (12,7);
  \draw[bkGray!90!black, line width=0.6pt] (4,5.5) -- (12,5.5);
  \bkthraspring{(7,5.5)}{9}\bkthraweight{(7,-3.5)}
  \filldraw[fill=bkYellowBg, draw=bkYellow!80!black, line width=0.5pt] (13,-12) rectangle (16.5,5.5);
  \foreach \k in {0,...,7} \draw[bkInk, line width=0.35pt] (13,{-10+\k*2}) -- (14.3,{-10+\k*2});
  \draw[bkRed, line width=0.5pt, densely dashed] (9.6,-3.5) -- (13,-3.5);
  % ② 表
  \draw[bkLine, line width=0.5pt] (29.5,4) rectangle (47,-6) (29.5,-1) -- (47,-1);
  \foreach \i in {1,2} \draw[bkLine, line width=0.5pt] ({29.5+\i*5.83},4) -- ({29.5+\i*5.83},-6);
  \foreach \v [count=\i from 0] in {0,1,2} \node[font=\footnotesize, inner sep=0pt] at ({32.4+\i*5.83},1.5) {\v};
  \foreach \v [count=\i from 0] in {0,1.6,2.9} \node[font=\footnotesize, inner sep=0pt] at ({32.4+\i*5.83},-3.5) {\v};
  % ③ 点
  \draw[-{Stealth[length=2.5pt]}, bkInk, line width=0.5pt] (58,-11) -- (75,-11);
  \draw[-{Stealth[length=2.5pt]}, bkInk, line width=0.5pt] (58,-11) -- (58,8);
  \foreach \x/\y in {0/0,1/1.6,2/2.9,3/4.6,4/5.9,5/7.6,6/8.9} \fill[bkInk] ({58+\x*2.3},{-11+\y*1.9}) circle (0.6);
  % ④ 線
  \draw[-{Stealth[length=2.5pt]}, bkInk, line width=0.5pt] (85.75,-11) -- (102.75,-11);
  \draw[-{Stealth[length=2.5pt]}, bkInk, line width=0.5pt] (85.75,-11) -- (85.75,8);
  \draw[bkBlue, line width=0.9pt] (85.75,-11) -- ({85.75+6.6*2.3},{-11+9.9*1.9});
  \foreach \x/\y in {0/0,1/1.6,2/2.9,3/4.6,4/5.9,5/7.6,6/8.9} \fill[bkInk] ({85.75+\x*2.3},{-11+\y*1.9}) circle (0.6);
  % ⑤ 式
  \node[font=\small, text=bkBlue!80!black] at (121.5,-1) {$y=1.5x$};
\end{tikzpicture}
\bkfigcaption{はかった\R{数}{かず}から\R{式}{しき}ができるまでの5つの\R{段}{だん}}
\end{bkfigure}

1つずつ、ばねの\R{実験}{じっ|けん}で\R{見}{み}ていきましょう。

\bksubsection{1}{はかる：\R{目}{め}もりを\R{読}{よ}む}

\noindent\begin{minipage}[c]{0.44\linewidth}\centering
\begin{tikzpicture}[x=1mm, y=1mm]
  % 天井と目もり板（1 cm を 2.5 mm でかく。0 はばねの上のはし）  % noruby
  \bkthraceil{(4,0)}{(54,0)}
  \filldraw[fill=bkYellowBg, draw=bkYellow!80!black, line width=0.6pt] (24,0) rectangle (34,-52.5);
  \foreach \k in {0,...,20} {\draw[bkInk, line width=0.35pt] (24,{-\k*2.5}) -- (25.4,{-\k*2.5});
    \draw[bkInk, line width=0.35pt] (32.6,{-\k*2.5}) -- (34,{-\k*2.5});}
  \foreach \k in {0,5,10,15,20} {\draw[bkInk, line width=0.7pt] (24,{-\k*2.5}) -- (26.4,{-\k*2.5});
    \draw[bkInk, line width=0.7pt] (31.6,{-\k*2.5}) -- (34,{-\k*2.5});}
  \foreach \k in {5,10,15,20} \node[font=\sffamily\footnotesize, inner sep=0pt] at (29,{-\k*2.5}) {\k};
  % おもり0個：長さ 10.0 cm  % noruby
  \bkthraspring{(14,0)}{25}
  \fill[bkRed] (14,-25) circle (0.9);
  \draw[bkRed, line width=0.6pt, densely dashed] (15.5,-25) -- (24,-25);
  \node[anchor=east, font=\sffamily\footnotesize, text=bkRed!80!black, inner sep=0pt] at (10.5,-25) {\qty{10.0}{cm}};
  % おもり3個：長さ 14.6 cm  % noruby
  \bkthraspring{(45,0)}{36.5}
  \bkthraweight{(45,-36.5)}\bkthraweight{(45,-43)}\bkthraweight{(45,-49.5)}
  \fill[bkRed] (41.5,-36.5) circle (0.9);
  \draw[bkRed, line width=0.6pt, densely dashed] (40,-36.5) -- (34,-36.5);
  \node[anchor=west, font=\sffamily\footnotesize, text=bkRed!80!black, inner sep=0pt] at (49.5,-36.5) {\qty{14.6}{cm}};
  \node[bk tick, anchor=north] at (14,-58) {0\R{個}{こ}};
  \node[bk tick, anchor=north] at (45,-58) {3\R{個}{こ}};
  \node[font=\sffamily\footnotesize, text=bkGray, anchor=north] at (29,-58) {（cm）};
\end{tikzpicture}
\end{minipage}\hfill
\begin{minipage}[c]{0.53\linewidth}
ソラさんは、ばねを\R{天井}{てん|じょう}の\R{板}{いた}につるし、\R{横}{よこ}に\R{目}{め}もり\R{板}{いた}を\R{立}{た}てました。\R{目}{め}もり\R{板}{いた}の0は、ばねの\R{上}{うえ}のはしにそろえてあります。だから、ばねの\R{下}{した}のはしの\R{目}{め}もりが、そのまま「ばねの\R{長}{なが}さ」です。

\R{目}{め}もり\R{板}{いた}の1\R{目}{め}もりは\qty{1}{cm}です。\R{目}{め}もりと\R{目}{め}もりのあいだは、1\R{目}{め}もりの10\R{分}{ぶん}の1まで、\R{目}{め}で\R{見}{み}つもって\R{読}{よ}みます（2-6）。おもりがないとき、ばねの\R{長}{なが}さは \qty{10.0}{cm} でした。

つぎに、1\R{個}{こ}\qty{20}{g}のおもりを1\R{個}{こ}ずつふやしてつるし、そのたびにばねがゆれなくなるのを\R{待}{ま}ってから、\R{目}{め}もりを\R{読}{よ}みました。3\R{個}{こ}のときは \qty{14.6}{cm} でした。
\end{minipage}

\par\vspace{2mm}
「\qty{10.0}{cm}」の「.0」は、\R{書}{か}かなくてもいいように\R{見}{み}えます。でも、「\R{小数第}{しょう|すう|だい}1\R{位}{い}まで\R{読}{よ}んだら、ちょうど0だった」という\R{大事}{だい|じ}な\R{記録}{き|ろく}です（2-6）。

\bksubsection{2}{\R{表}{ひょう}：はかった\R{数}{かず}を\R{並}{なら}べる}

はかった\R{数}{かず}は、ノートにばらばらに\R{書}{か}かないで、\R{表}{ひょう}に\R{並}{なら}べます。\R{表}{ひょう}にすると、\R{変}{か}わり\R{方}{かた}が\R{横}{よこ}に\R{並}{なら}んで\R{見}{み}えるからです。

\begin{bkcenter}
\begin{bktable}{colspec={Q[l,m,wd=30mm] *{7}{Q[c,m,wd=9mm]}}, rowsep=1.5pt, cells={font=\small}}
おもりの\R{数}{かず} $x$（\R{個}{こ}） & 0 & 1 & 2 & 3 & 4 & 5 & 6 \\
ばねの\R{長}{なが}さ（cm） & 10.0 & 11.6 & 12.9 & 14.6 & 15.9 & 17.6 & 18.9 \\
\SetCell{bg=bkYellow!30} のび $y$（cm） & 0 & 1.6 & 2.9 & 4.6 & 5.9 & 7.6 & 8.9 \\
\end{bktable}
\bktabcaption{ソラさんがはかったばねの\R{長}{なが}さと、のび}
\end{bkcenter}

いちばん\R{下}{した}の\R{行}{ぎょう}の「のび」は、ばねの\R{長}{なが}さから、おもりがないときの\R{長}{なが}さ\qty{10.0}{cm}をひいた\R{数}{かず}です。「あと－まえ」で、\R{変}{か}わった\R{量}{りょう}を\R{求}{もと}めました（2-3）。たとえば3\R{個}{こ}のときは、$14.6-10.0=4.6$ で、\qty{4.6}{cm}のびています。

\R{表}{ひょう}を\R{横}{よこ}に\R{見}{み}ると、おもりが1\R{個}{こ}ふえるごとに、のびは1.6、1.3、1.7、1.3、1.7、1.3 とふえています（となりどうしのひき\R{算}{ざん}）。だいたい1.5くらいずつ、でも、ぴったり\R{同}{おな}じではありません。これは、あとでもう\R{一度}{いち|ど}\R{考}{かんが}えます。

\bksubsection{3}{\R{点}{てん}：\R{表}{ひょう}の1\R{列}{れつ}を、グラフの1つの\R{点}{てん}にする}

\noindent\begin{minipage}[c]{0.5\linewidth}\centering
\begin{tikzpicture}
\begin{axis}[bkplot, width=62mm, height=56mm, xmin=0, xmax=6.8, ymin=0, ymax=10.5, xtick={0,...,6}, ytick={0,2,...,10},
  xlabel={$x$（\R{個}{こ}）}, ylabel={$y$（cm）}]
  \addplot[bkBlue, dashed, line width=0.8pt] coordinates {(3,0) (3,4.6) (0,4.6)};
  \addplot[only marks, mark=*, mark size=1.8pt, bkInk] coordinates {(0,0) (1,1.6) (2,2.9) (4,5.9) (5,7.6) (6,8.9)};
  \addplot[only marks, mark=*, mark size=2.4pt, bkRed] coordinates {(3,4.6)};
  \node[font=\sffamily\footnotesize, anchor=south east, inner sep=1pt, text=bkRed!80!black] at (axis cs:2.9,4.9) {3\R{個}{こ}で4.6};
\end{axis}
\end{tikzpicture}
\end{minipage}\hfill
\begin{minipage}[c]{0.47\linewidth}
\R{横}{よこ}の\R{軸}{じく}におもりの\R{数}{かず} $x$、たての\R{軸}{じく}にのび $y$ をとります。\R{表}{ひょう}のたての1\R{列}{れつ}（$x$ と $y$ の1\R{組}{くみ}）を、1つの\R{点}{てん}にします。

\R{赤}{あか}い\R{点}{てん}は、3\R{個}{こ}で\qty{4.6}{cm}のびた、という1\R{組}{くみ}です。0から\R{右}{みぎ}へ3、そこから\R{上}{うえ}へ4.6\R{進}{すす}んだ\R{所}{ところ}にかきます（2-15 の\R{点}{てん}のかき\R{方}{かた}。3-2 でくわしく\R{学}{まな}びます）。
\end{minipage}

\par\vspace{2mm}
7\R{列}{れつ}の\R{表}{ひょう}は、7つの\R{点}{てん}になりました。\R{数}{かず}の\R{並}{なら}びでは\R{見}{み}えにくかったことが、\R{点}{てん}にすると\R{一目}{ひと|め}で\R{見}{み}えます。\R{点}{てん}は、\R{左下}{ひだり|した}から\R{右上}{みぎ|うえ}へ、ほぼまっすぐに\R{並}{なら}んでいます。

\bksubsection{4}{\R{線}{せん}：\R{点}{てん}の\R{並}{なら}びに\R{合}{あ}う\R{線}{せん}を\R{引}{ひ}く}

\noindent\begin{minipage}[c]{0.47\linewidth}
\R{点}{てん}がまっすぐに\R{並}{なら}んでいるので、\R{透明}{とう|めい}な\R{定規}{じょう|ぎ}を\R{当}{あ}てて、\R{点}{てん}のなるべく\R{近}{ちか}くを\R{通}{とお}る\R{直線}{ちょく|せん}を1\R{本}{ぽん}\R{引}{ひ}きます。おもりが0\R{個}{こ}ならのびも0なので、\R{線}{せん}は\R{原点}{げん|てん}（$x=0$、$y=0$ の\R{点}{てん}）から\R{引}{ひ}きます。

\R{点}{てん}どうしを\R{順}{じゅん}に\R{結}{むす}んだ、ぎざぎざのおれ\R{線}{せん}にはしません。\R{点}{てん}のでこぼこは、ばねのきまりではなく、はかるときのずれ（2-6）だと\R{考}{かんが}えるからです（3-4 でくわしく）。
\end{minipage}\hfill
\begin{minipage}[c]{0.5\linewidth}\centering
\begin{tikzpicture}
\begin{axis}[bkplot, width=62mm, height=56mm, xmin=0, xmax=6.8, ymin=0, ymax=10.5, xtick={0,...,6}, ytick={0,2,...,10},
  xlabel={$x$（\R{個}{こ}）}, ylabel={$y$（cm）}]
  \addplot[bkcurve, domain=0:6.6] {1.5*x};
  \addplot[only marks, mark=*, mark size=1.8pt, bkInk] coordinates {(0,0) (1,1.6) (2,2.9) (3,4.6) (4,5.9) (5,7.6) (6,8.9)};
  \addplot[bkGreen!70!black, line width=1pt] coordinates {(2,3) (4,3) (4,6)};
  \node[font=\sffamily\footnotesize, text=bkGreen!50!black, anchor=north] at (axis cs:3,2.85) {\R{右}{みぎ}へ2};
  \node[font=\sffamily\footnotesize, text=bkGreen!50!black, anchor=west] at (axis cs:4.1,4.3) {\R{上}{うえ}へ3};
  \node[font=\sffamily\footnotesize, text=bkBlue!80!black, anchor=east] at (axis cs:5.2,9.4) {\R{引}{ひ}いた\R{直線}{ちょく|せん}};
\end{axis}
\end{tikzpicture}
\end{minipage}

\bksubsection{5}{\R{式}{しき}：\R{線}{せん}のきまりを\R{文字}{も|じ}で\R{書}{か}く}

\R{引}{ひ}いた\R{直線}{ちょく|せん}を\R{読}{よ}みます。\R{緑}{みどり}の\R{階段}{かい|だん}のように、\R{直線}{ちょく|せん}は\R{右}{みぎ}へ2\R{進}{すす}むと、\R{上}{うえ}へ3\R{上}{あ}がります。だから、\R{右}{みぎ}へ1\R{進}{すす}むと $3\div2=1.5$ \R{上}{あ}がります。おもり1\R{個}{こ}あたり\qty{1.5}{cm}のびる、ということです。

1\R{個}{こ}あたり\qty{1.5}{cm}なので、$x$ \R{個}{こ}なら、\qty{1.5}{cm}の $x$ こ\R{分}{ぶん}で $1.5\times x$（cm）のびます。$x=0$ なら $y=0$ で、\R{原点}{げん|てん}も\R{通}{とお}ります。\R{文字}{も|じ}の\R{式}{しき}では×をはぶいて（2-12）、
\[ y=1.5x \]
と\R{書}{か}けます。2-15 で\R{学}{まな}んだ\R{比例}{ひ|れい}の\R{形}{かたち} $y=ax$ で、1.5が\R{比例定数}{ひ|れい|てい|すう}です。

\bkhatenaref の\R{答}{こた}え：おもりを1\R{個}{こ}ふやすごとに、のびは\R{約}{やく}\qty{1.5}{cm}ずつふえます。\R{表}{ひょう}とグラフから、のびを\R{計算}{けい|さん}する\R{式}{しき} $y=1.5x$ を\R{作}{つく}れました。

\begin{taisetsu}
\R{式}{しき}ができると、はかっていない\R{数}{かず}も\R{計算}{けい|さん}できる。\par
おもり8\R{個}{こ}なら、$y=1.5\times8=12$ で、\R{約}{やく}\qty{12}{cm}のびる。
\end{taisetsu}

ただし、ばねは\R{引}{ひ}っぱりすぎると、のびきってもとにもどらなくなり、\R{式}{しき}が\R{使}{つか}えなくなります。\R{式}{しき}が\R{使}{つか}える $x$ の\R{範囲}{はん|い}は 3-1 で、\R{線}{せん}をはかった\R{所}{ところ}より\R{先}{さき}へのばす\R{危}{あぶ}なさは 3-4 で\R{学}{まな}びます。

\noindent\begin{minipage}[c]{0.5\linewidth}\centering
\begin{tikzpicture}
\begin{axis}[bkplot, width=60mm, height=56mm, xmin=0, xmax=6.8, ymin=0, ymax=21, xtick={0,...,6}, ytick={0,5,...,20},
  xlabel={$x$（\R{個}{こ}）}, ylabel={\R{長}{なが}さ（cm）}]
  \addplot[bkcurve2, domain=0:6.6] {10+1.5*x};
  \addplot[only marks, mark=*, mark size=1.8pt, bkInk] coordinates {(0,10.0) (1,11.6) (2,12.9) (3,14.6) (4,15.9) (5,17.6) (6,18.9)};
  \node[font=\sffamily\footnotesize, text=bkOrange!80!black, anchor=north west] at (axis cs:0.15,9.4) {$x=0$ で10};
\end{axis}
\end{tikzpicture}
\end{minipage}\hfill
\begin{minipage}[c]{0.47\linewidth}
\textbf{\R{長}{なが}さでグラフにすると？}\par
のびのかわりに、ばねの\R{長}{なが}さ（\R{表}{ひょう}のまん\R{中}{なか}の\R{行}{ぎょう}）で\R{点}{てん}をうつと、\R{左}{ひだり}のグラフになります。\R{点}{てん}はやはりまっすぐに\R{並}{なら}びますが、\R{線}{せん}は\R{原点}{げん|てん}を\R{通}{とお}らず、10の\R{所}{ところ}から\R{始}{はじ}まります。\R{長}{なが}さは $10+1.5x$ で、\R{比例}{ひ|れい}ではありません。このような\R{直線}{ちょく|せん}の\R{式}{しき}は、3-5 でくわしく\R{学}{まな}びます。
\end{minipage}

%% ---------------- 点は線にぴったりのらない ----------------
\bksection{me}{\R{点}{てん}は、\R{線}{せん}の\R{上}{うえ}にぴったりのらない}

\R{表}{ひょう}の\R{数}{かず}で、「のび÷おもりの\R{数}{かず}」を\R{計算}{けい|さん}してみましょう。おもり1\R{個}{こ}あたりののびです（2-3 の「1あたりの\R{量}{りょう}」）。

\begin{bkcenter}
\begin{bktable}{colspec={Q[l,m,wd=36mm] *{6}{Q[c,m,wd=11mm]}}, rowsep=1.5pt, cells={font=\small}}
おもりの\R{数}{かず} $x$（\R{個}{こ}） & 1 & 2 & 3 & 4 & 5 & 6 \\
のび $y$（cm） & 1.6 & 2.9 & 4.6 & 5.9 & 7.6 & 8.9 \\
\SetCell{bg=bkYellow!30} $y\div x$（cm/\R{個}{こ}） & 1.6 & 1.45 & 1.53… & 1.475 & 1.52 & 1.48… \\
\end{bktable}
\bktabcaption{1\R{個}{こ}あたりののびは、どれも1.5の\R{近}{ちか}く。でも、ぴったり1.5ではない。}
\end{bkcenter}

どれも1.5の\R{近}{ちか}くですが、1つも1.5ぴったりではありません。だから、\R{点}{てん}は\R{直線}{ちょく|せん} $y=1.5x$ の\R{少}{すこ}し\R{上}{うえ}や\R{少}{すこ}し\R{下}{した}にあり、\R{線}{せん}の\R{上}{うえ}にぴったりのってはいません。

\R{理由}{り|ゆう}は、はかるときのずれです（2-6）。\R{目}{め}もりの10\R{分}{ぶん}の1は\R{目}{め}で\R{見}{み}つもるので、\R{人}{ひと}によって\R{少}{すこ}しちがいます。ばねはなかなかゆれが\R{止}{と}まりません。おもりも、どれもぴったり\qty{20}{g}とはかぎりません。どんなにていねいにはかっても、ずれは0にはなりません。

\noindent\begin{minipage}[c]{0.4\linewidth}\centering
\begin{tikzpicture}
\begin{axis}[bkplot, width=50mm, height=50mm, xmin=0, xmax=6.8, ymin=0, ymax=10.5, xtick={0,2,4,6}, ytick={0,2,...,10},
  xlabel={$x$}, ylabel={$y$}]
  \addplot[bkOrange, line width=0.9pt, domain=0:6.6] {1.45*x};
  \addplot[bkBlue, line width=0.9pt, domain=0:6.6] {1.5*x};
  \addplot[bkGreen!80!black, line width=0.9pt, domain=0:6.6] {1.55*x};
  \addplot[only marks, mark=*, mark size=1.5pt, bkInk] coordinates {(0,0) (1,1.6) (2,2.9) (3,4.6) (4,5.9) (5,7.6) (6,8.9)};
  \draw[bkRed, line width=0.8pt] (axis cs:3.5,4) rectangle (axis cs:6.5,10.2);
\end{axis}
\end{tikzpicture}
\end{minipage}\hfill
\begin{minipage}[c]{0.58\linewidth}\centering
\begin{tikzpicture}
\begin{axis}[bkplot, axis lines=box, axis line style={-, line width=0.8pt, bkRed}, width=70mm, height=60mm, xmin=3.5, xmax=6.5, ymin=4, ymax=10.2,
  xtick={4,5,6}, ytick={4,5,...,10}, xlabel={}, ylabel={}]
  \addplot[bkOrange, line width=1.2pt, domain=3.5:6.5] {1.45*x};
  \addplot[bkBlue, line width=1.2pt, domain=3.5:6.5] {1.5*x};
  \addplot[bkGreen!80!black, line width=1.2pt, domain=3.5:6.5] {1.55*x};
  \addplot[only marks, mark=*, mark size=2pt, bkInk] coordinates {(4,5.9) (5,7.6) (6,8.9)};
  \draw[bkGreen!80!black, line width=1.2pt] (axis cs:3.65,9.75) -- (axis cs:3.95,9.75);
  \node[anchor=west, font=\footnotesize, text=bkGreen!50!black] at (axis cs:3.97,9.75) {$y=1.55x$};
  \draw[bkBlue, line width=1.2pt] (axis cs:3.65,9.05) -- (axis cs:3.95,9.05);
  \node[anchor=west, font=\footnotesize, text=bkBlue!80!black] at (axis cs:3.97,9.05) {$y=1.5x$};
  \draw[bkOrange, line width=1.2pt] (axis cs:3.65,8.35) -- (axis cs:3.95,8.35);
  \node[anchor=west, font=\footnotesize, text=bkOrange!80!black] at (axis cs:3.97,8.35) {$y=1.45x$};
\end{axis}
\end{tikzpicture}
\end{minipage}
\bkfigcaption{\R{左}{ひだり}のグラフの\R{赤}{あか}い\R{四角}{し|かく}の\R{中}{なか}を、\R{右}{みぎ}で\R{大}{おお}きくした。3\R{本}{ぼん}の\R{直線}{ちょく|せん}は、どれも\R{点}{てん}の\R{近}{ちか}くを\R{通}{とお}る。}

\R{右}{みぎ}のグラフでは、3\R{本}{ぼん}の\R{直線}{ちょく|せん}を\R{引}{ひ}きました。どの\R{線}{せん}も、\R{点}{てん}の\R{近}{ちか}くを\R{通}{とお}ります。でも、どの\R{線}{せん}にも、\R{上}{うえ}にはずれた\R{点}{てん}と\R{下}{した}にはずれた\R{点}{てん}があります。\R{友}{とも}だちと\R{同}{おな}じ\R{表}{ひょう}からグラフをかいて、\R{定規}{じょう|ぎ}で\R{線}{せん}を\R{引}{ひ}くと、きっと\R{少}{すこ}しずつちがう\R{線}{せん}になるでしょう。

\begin{hatena}
\textbf{この\R{巻}{かん}からの\R{宿題}{しゅく|だい}}\par\noindent
ずれた\R{点}{てん}に、どの\R{直線}{ちょく|せん}を\R{引}{ひ}けばいい？\quad \R{人}{ひと}によって\R{答}{こた}えが\R{変}{か}わらない\R{決}{き}め\R{方}{かた}はある？
\end{hatena}

「\R{目}{め}で\R{見}{み}て、だいたい」ではなく、だれがやっても\R{同}{おな}じ\R{線}{せん}になる\R{決}{き}め\R{方}{かた}があると\R{便利}{べん|り}です。そのためには、まず「ずれの\R{大}{おお}きさ」を1つの\R{数}{かず}で\R{言}{い}えるようになる\R{必要}{ひつ|よう}があります（3-13・3-14）。そして、ずれ\R{全体}{ぜん|たい}がいちばん\R{小}{ちい}さくなる\R{直線}{ちょく|せん}を、\R{計算}{けい|さん}で\R{決}{き}める\R{方法}{ほう|ほう}を、この\R{巻}{かん}の\R{終}{お}わりの 3-16 で\R{学}{まな}びます。この\R{宿題}{しゅく|だい}を、\R{覚}{おぼ}えておいてください。

\begin{yattemiyou}[\R{輪}{わ}ゴムでもやってみよう]
\begin{enumerate}[label={\bkcircled[bkGreen]{\arabic*}}, leftmargin=2.4\zw, itemsep=1pt, topsep=2pt]
\item \R{輪}{わ}ゴムをフックにかけ、\R{下}{した}に\R{小}{ちい}さなポリぶくろをつるそう。\R{横}{よこ}にものさしを\R{立}{た}てて、ふくろの\R{上}{うえ}のはしの\R{目}{め}もりを\R{読}{よ}もう。
\item ふくろに10\R{円玉}{えん|だま}を5まいずつ\R{入}{い}れて、そのたびに\R{目}{め}もりを\R{読}{よ}もう（10\R{円玉}{えん|だま}1まいは\R{約}{やく}\qty{4.5}{g}）。
\item のびを\R{表}{ひょう}にして、\R{方眼}{ほう|がん}のノートに\R{点}{てん}をうとう。\R{横}{よこ}の\R{軸}{じく}は10\R{円玉}{えん|だま}のまい\R{数}{すう}、たての\R{軸}{じく}はのび（cm）。
\item \R{点}{てん}はまっすぐに\R{並}{なら}んだかな？\quad ばねのグラフとくらべよう。
\end{enumerate}
\R{輪}{わ}ゴムは、のびるにつれてのび\R{方}{かた}が\R{変}{か}わり、\R{点}{てん}が\R{曲}{ま}がって\R{並}{なら}ぶことがあります。まっすぐでないグラフのきまりは、\R{第}{だい}2\R{部}{ぶ}（3-7〜3-10）で\R{学}{まな}びます。
\end{yattemiyou}

\begin{column}{なぜ「\R{関数}{かん|すう}とデータ」なの？}
\R{物理}{ぶつ|り}のきまりの\R{多}{おお}くは、「この\R{量}{りょう}が\R{決}{き}まると、あの\R{量}{りょう}が\R{決}{き}まる」という\R{形}{かたち}をしています。おもりの\R{数}{かず}が\R{決}{き}まると、ばねののびが\R{決}{き}まる。\R{時刻}{じ|こく}が\R{決}{き}まると、ボールの\R{高}{たか}さが\R{決}{き}まる。このような\R{関係}{かん|けい}を、\R{数学}{すう|がく}では\textbf{\R{関数}{かん|すう}}とよびます（3-3）。\bkidx{かんすう}{\R{関数}{かん|すう}}

\R{関数}{かん|すう}は、\R{頭}{あたま}の\R{中}{なか}だけで\R{見}{み}つかるものではありません。はかった\R{数}{かず}（データ）を\R{表}{ひょう}とグラフにして、はじめて\R{見}{み}えてきます。\R{関数}{かん|すう}とデータは、きまりを\R{見}{み}つけるための\R{車}{くるま}の\R{両輪}{りょう|りん}です。\R{技術者}{ぎ|じゅつ|しゃ}は、\R{実験}{じっ|けん}のデータから\R{式}{しき}を\R{作}{つく}り、その\R{式}{しき}で\R{機械}{き|かい}を\R{設計}{せっ|けい}します。
\end{column}


\begin{yattemiyou}[\R{表計算}{ひょう|けい|さん}で、5つの\R{段}{だん}をやってみよう]
\noindent\begin{minipage}[c]{0.56\linewidth}
\begin{enumerate}[label={\bkcircled[bkGreen]{\arabic*}}, leftmargin=2.4\zw, itemsep=1pt, topsep=0pt]
\item A\R{列}{れつ}におもりの\R{数}{かず}、B\R{列}{れつ}にばねの\R{長}{なが}さを\R{入力}{にゅう|りょく}しよう。
\item C2に「=B2-B\$2」と\R{入力}{にゅう|りょく}して、C8までコピーしよう（\$ は 2-18）。
\item A\R{列}{れつ}とC\R{列}{れつ}を\R{選}{えら}んで、\R{点}{てん}だけのグラフ（\R{散布図}{さん|ぷ|ず}）を\R{作}{つく}ろう。
\item グラフに「\R{近似曲線}{きん|じ|きょく|せん}（\R{線形}{せん|けい}）」を\R{追加}{つい|か}して、\R{式}{しき}を\R{表示}{ひょう|じ}させよう。
\end{enumerate}
\end{minipage}\hfill
\begin{minipage}[c]{0.4\linewidth}\centering
\bksheet{3}{11}{%
1 & $x$ & \R{長}{なが}さ & のび \\
2 & 0 & 10.0 & 0.0 \\
3 & 1 & 11.6 & 1.6 \\
4 & 2 & 12.9 & 2.9 \\
5 & 3 & 14.6 & 4.6 \\
6 & 4 & 15.9 & 5.9 \\
7 & 5 & 17.6 & 7.6 \\
8 & 6 & 18.9 & 8.9 \\}
\end{minipage}
\par\vspace{4pt}
④では、$y=1.4893x+0.0321$ のような\R{式}{しき}が\R{出}{で}ます（\R{表計算}{ひょう|けい|さん}のソフトによって、\R{書}{か}き\R{方}{かた}は\R{少}{すこ}しちがいます）。1.4893は1.5にとても\R{近}{ちか}く、0.0321は0にとても\R{近}{ちか}い\R{数}{かず}です。\R{表計算}{ひょう|けい|さん}は、\R{定規}{じょう|ぎ}を\R{使}{つか}わずに、\R{計算}{けい|さん}で1\R{本}{ぽん}の\R{直線}{ちょく|せん}を\R{決}{き}めてくれました。どうやって\R{決}{き}めているのかは、3-16 でわかります。
\end{yattemiyou}

ただし、\R{表計算}{ひょう|けい|さん}が\R{出}{だ}した\R{式}{しき}を、そのまま\R{信}{しん}じてはいけません。\R{数}{かず}を1つ\R{入力}{にゅう|りょく}しまちがえても、\R{表計算}{ひょう|けい|さん}はだまって\R{式}{しき}を\R{出}{だ}します。だから、\R{式}{しき}を\R{使}{つか}う\R{前}{まえ}に、かならずグラフの\R{点}{てん}を\R{自分}{じ|ぶん}の\R{目}{め}で\R{見}{み}て、1つだけ\R{遠}{とお}くはなれた\R{点}{てん}がないか、\R{線}{せん}が\R{点}{てん}の\R{並}{なら}びに\R{合}{あ}っているかを\R{確}{たし}かめます。\R{機械}{き|かい}に\R{計算}{けい|さん}をまかせても、\R{確}{たし}かめるのは\R{人}{ひと}の\R{仕事}{し|ごと}です。

\bkfukidashi{sora}{\R{線}{せん}を\R{引}{ひ}くたびに、ちょっとずつ\R{傾}{かたむ}きが\R{変}{か}わっちゃう。\R{計算}{けい|さん}で1つに\R{決}{き}まるなら、すっきりするね。}

\begin{butsuri}
\par\addvspace{3pt}\textbf{はかりの\R{目}{め}もりを\R{決}{き}める：\R{校正}{こう|せい}}\par\noindent
\begin{minipage}[c]{0.62\linewidth}
キッチンスケールや\R{体重計}{たい|じゅう|けい}の\R{中}{なか}には、\R{重}{おも}さがかかると\R{少}{すこ}しだけ\R{曲}{ま}がる\R{金属}{きん|ぞく}の\R{部品}{ぶ|ひん}（ロードセル）が\R{入}{はい}っています。\R{曲}{ま}がると、\R{部品}{ぶ|ひん}から\R{出}{で}る\R{電気}{でん|き}の\R{信号}{しん|ごう}（\R{電圧}{でん|あつ}）が\R{少}{すこ}し\R{変}{か}わります。
\end{minipage}\hfill
\begin{minipage}[c]{0.37\linewidth}\centering
\begin{tikzpicture}[x=1mm, y=1mm]
  % 台・ロードセル（左を固定した棒）・のせ台・分銅・表示
  \fill[bkGray!40, rounded corners=1mm] (0,0) rectangle (40,3);
  \fill[bkGray!70] (3,3) rectangle (9,9);
  \fill[bkBlue!45] (9,6.2) rectangle (31,8.8);
  \draw[bkRed, line width=0.7pt, densely dashed] (9,6.2) .. controls (20,6) and (26,5.2) .. (31,4.6);
  \fill[bkGray!70] (27,8.8) rectangle (30,12);
  \fill[bkGray!55, rounded corners=0.6mm] (12,12) rectangle (38,14);
  \filldraw[fill=bkGray!45, draw=bkGray!90!black, line width=0.6pt, rounded corners=0.8mm] (20,14) rectangle (30,22);
  \node[font=\sffamily\footnotesize, text=bkInk] at (25,18) {1\,kg};
  \node[anchor=north, font=\sffamily\footnotesize, text=bkBlue!80!black, inner sep=0pt] at (20,-1.5) {ロードセル};
\end{tikzpicture}
\end{minipage}
\par\vspace{2pt}
\R{工場}{こう|じょう}では、\R{重}{おも}さが\R{正確}{せい|かく}にわかっているおもり（\R{分銅}{ふん|どう}）をのせて\R{信号}{しん|ごう}をはかり、\R{表}{ひょう}とグラフから\R{直線}{ちょく|せん}の\R{式}{しき}を\R{作}{つく}ります。1\,kg ごとに\R{信号}{しん|ごう}が \qty{2.5}{mV}（\qty{1}{V}の1000\R{分}{ぶん}の1が\qty{1}{mV}）ふえるなら、\R{信号}{しん|ごう}が \qty{3.0}{mV} のときは $3.0\div2.5=1.2$ で、\qty{1.2}{kg}です。\R{計器}{けい|き}の\R{目}{め}もりをこのように\R{決}{き}めることを、\textbf{\R{校正}{こう|せい}}といいます。

\par\addvspace{3pt}\textbf{データロガー：はかる・\R{表}{ひょう}にするを\R{自動}{じ|どう}で}\par\noindent
\R{工場}{こう|じょう}や\R{気象台}{き|しょう|だい}では、センサーが1\R{秒}{びょう}ごとや10\R{分}{ぷん}ごとに\R{温度}{おん|ど}や\R{電圧}{でん|あつ}をはかり、\R{自動}{じ|どう}で\R{表}{ひょう}にしていく\R{機械}{き|かい}（データロガー）が\R{使}{つか}われています。1\R{日}{にち}で\R{何万行}{なん|まん|ぎょう}にもなる\R{表}{ひょう}は、\R{人}{ひと}の\R{目}{め}では\R{読}{よ}みきれません。グラフにすれば、どこで\R{何}{なに}が\R{起}{お}きたかが\R{一目}{ひと|め}でわかります（3-4）。

\par\addvspace{3pt}\textbf{\R{大学}{だい|がく}への\R{種}{たね}}\par\noindent
\R{実験}{じっ|けん}のデータから\R{機械}{き|かい}の\R{式}{しき}を\R{作}{つく}ることは、\R{大学}{だい|がく}の\R{制御}{せい|ぎょ}の\R{勉強}{べん|きょう}でも\R{大事}{だい|じ}な\R{仕事}{し|ごと}です。\R{式}{しき}があれば、\R{機械}{き|かい}を\R{動}{うご}かす\R{前}{まえ}に、どう\R{動}{うご}くかを\R{計算}{けい|さん}で\R{確}{たし}かめられるからです。
\end{butsuri}









\clearpage
%% ---------------- 自由研究のヒント ----------------
\bksection{te}{はかって、きまりを\R{見}{み}つけよう}

ばねのほかにも、\R{身}{み}のまわりには、はかって\R{表}{ひょう}とグラフにすると、きまりが\R{見}{み}えてくるものがたくさんあります。\R{夏休}{なつ|やす}みの\R{自由研究}{じ|ゆう|けん|きゅう}にもぴったりです。グラフがどんな\R{形}{かたち}になるか、\R{先}{さき}に\R{予想}{よ|そう}してから、はかってみましょう。

\begin{bkcenter}
\begin{bktable}{colspec={Q[l,m,wd=50mm] Q[l,m,wd=21mm] Q[l,m,wd=19mm] Q[c,m,wd=11mm] Q[c,m,wd=9mm]}, rowsep=0.5pt, cells={font=\small}, row{1}={font=\sffamily\bfseries\small}}
はかること & \R{横}{よこ}の\R{軸}{じく} & たての\R{軸}{じく} & \R{形}{かたち} & \R{章}{しょう} \\
\R{同}{おな}じ\R{速}{はや}さで\R{歩}{ある}いた\R{歩数}{ほ|すう}（スマホの\R{歩数計}{ほ|すう|けい}） & \R{時間}{じ|かん}（\R{分}{ふん}） & \R{歩数}{ほ|すう} & ア & 3-5 \\
\R{振}{ふ}り\R{子}{こ}が10\R{回}{かい}ゆれる\R{時間}{じ|かん} & \R{長}{なが}さ（cm） & \R{時間}{じ|かん}（\R{秒}{びょう}） & イ & 3-11 \\
\R{懐中電灯}{かい|ちゅう|でん|とう}からはなれたときの\R{明}{あか}るさ（スマホの\R{照度}{しょう|ど}センサー） & \R{距離}{きょ|り}（cm） & \R{明}{あか}るさ（lx） & ウ & 3-7 \\
\R{坂}{さか}を\R{転}{ころ}がるボールの\R{進}{すす}んだ\R{距離}{きょ|り}（\R{動画}{どう|が}で\R{調}{しら}べる） & \R{時間}{じ|かん}（\R{秒}{びょう}） & \R{距離}{きょ|り}（cm） & エ & 3-9 \\
\end{bktable}
\bktabcaption{はかって、グラフにしてみよう（「グラフの\R{形}{かたち}」は、\R{下}{した}の\R{図}{ず}のア〜エ）}
\end{bkcenter}

\begin{bkfigure}
\begin{tikzpicture}[x=1mm, y=1mm]
  \foreach \cx/\lab/\txt in {0/ア/{まっすぐ}, 34/イ/{だんだん\\ゆるやかに\R{上}{あ}がる}, 68/ウ/{だんだん\\ゆるやかに\R{下}{さ}がる}, 102/エ/{だんだん\\\R{急}{きゅう}に\R{上}{あ}がる}} {
    \draw[-{Stealth[length=3pt]}, bkInk, line width=0.6pt] (\cx,0) -- (\cx+28,0);
    \draw[-{Stealth[length=3pt]}, bkInk, line width=0.6pt] (\cx,0) -- (\cx,20);
    \node[anchor=north, font=\sffamily\bfseries\small, text=bkDeep] at (\cx+14,-2) {\lab};
    \node[anchor=north, font=\sffamily\footnotesize, align=center, text=bkInk] at (\cx+14,-7.5) {\txt};}
  \draw[bkBlue, line width=1.3pt] (0,0) -- (24,17);
  \draw[bkBlue, line width=1.3pt, domain=0:24, samples=40] plot ({34+\x}, {3.5*sqrt(\x)});
  \draw[bkBlue, line width=1.3pt, domain=1.2:24, samples=40] plot ({68+\x}, {20/\x});
  \draw[bkBlue, line width=1.3pt, domain=0:24, samples=40] plot ({102+\x}, {\x*\x/34});
\end{tikzpicture}
\bkfigcaption{グラフの\R{形}{かたち}の\R{例}{れい}。\R{形}{かたち}がわかると、\R{式}{しき}のなかまが\R{予想}{よ|そう}できる。}
\end{bkfigure}

まっすぐなグラフ（ア）なら、ばねと\R{同}{おな}じように\R{傾}{かたむ}きを\R{読}{よ}めば\R{式}{しき}が\R{作}{つく}れます。\R{曲}{ま}がったグラフ（イ・ウ・エ）は、どんな\R{式}{しき}になるのでしょう。それを\R{見分}{み|わ}ける\R{道具}{どう|ぐ}を、この\R{巻}{かん}で1つずつ\R{手}{て}に\R{入}{い}れていきます。\R{表}{ひょう}の\R{右}{みぎ}のはしの\R{章}{しょう}を\R{読}{よ}んだら、もう\R{一度}{いち|ど}このページにもどってきてください。

\begin{taisetsu}
はかった\R{数}{かず}は、\R{消}{け}さずにそのまま\R{記録}{き|ろく}しよう（へんだと\R{思}{おも}った\R{数}{かず}も）。
\end{taisetsu}

\clearpage
%% ---------------- この巻が通る路線と駅 ----------------
\bksection{bkthraroute}{この\R{巻}{かん}が\R{通}{とお}る\R{路線}{ろ|せん}と\R{駅}{えき}}

この\R{本}{ほん}の\R{章}{しょう}は、7\R{本}{ほん}の\R{路線}{ろ|せん}の\R{駅}{えき}です。\R{同}{おな}じ\R{路線}{ろ|せん}の\R{駅}{えき}は、1つの\R{考}{かんが}え\R{方}{かた}でつながっています。\R{第}{だい}3\R{巻}{かん}を\R{通}{とお}るのは、\R{次}{つぎ}の5\R{本}{ほん}です。

\begin{bkfigure}
\begin{tikzpicture}[x=1mm, y=1mm]
  \bkthraroute{48}{1}{2-3}{{\bfseries 3-5}・3-7・3-8・3-9・3-10}{5-1}
  \bkthraroute{36}{3}{2-7}{{\bfseries 3-11}・3-12}{5-14}
  \bkthraroute{24}{5}{2-18}{3-3・{\bfseries 3-6}}{7-1}
  \bkthraroute{12}{6}{2-2}{{\bfseries 3-2}}{4-10}
  \bkthraroute{0}{7}{2-6}{3-1・3-4・3-13・{\bfseries 3-14}・3-15・3-16・3-17}{9-4}
  \node[font=\sffamily\bfseries\small, text=bkGray] at (24,58) {\R{前}{まえ}の\R{巻}{かん}};
  \node[font=\sffamily\bfseries\small, text=bkDeep] at (74,58) {この\R{巻}{かん}（\R{第}{だい}3\R{巻}{かん}）};
  \node[font=\sffamily\bfseries\small, text=bkGray] at (123,58) {この\R{先}{さき}};
\end{tikzpicture}
\bkfigcaption{\R{第}{だい}3\R{巻}{かん}を\R{通}{とお}る5\R{本}{ほん}の\R{路線}{ろ|せん}。\R{四角}{し|かく}の\R{中}{なか}がこの\R{巻}{かん}の\R{駅}{えき}（\R{太字}{ふと|じ}は\R{大}{おお}きな\R{駅}{えき}）。}
\end{bkfigure}

\begin{bkcenter}
\begin{bktable}{colspec={Q[c,m,wd=7mm] Q[l,m,wd=16mm] Q[l,m,wd=92mm]}, rowsep=0.8pt, cells={font=\small}, row{1}={font=\sffamily\bfseries\small}}
 & \R{路線}{ろ|せん} & \R{大}{おお}きな\R{駅}{えき}のつながり \\
\bkrouteicon{1} & \R{変化}{へん|か} & \R{変化量}{へん|か|りょう}（2-3）→ \R{傾}{かたむ}き（3-5）→ \R{微分}{び|ぶん}（5-1） \\
\bkrouteicon{3} & \R{倍々}{ばい|ばい} & \R{累乗}{るい|じょう}（2-7）→ \R{指数関数}{し|すう|かん|すう}（3-11）→ $e$（5-14） \\
\bkrouteicon{5} & \R{表}{ひょう} & \R{表計算}{ひょう|けい|さん}（2-18）→ \R{連立方程式}{れん|りつ|ほう|てい|しき}（3-6）→ \R{行列}{ぎょう|れつ}（7-1） \\
\bkrouteicon{6} & \R{矢印}{や|じるし} & マイナスの\R{向}{む}き（2-2）→ \R{座標}{ざ|ひょう}（3-2）→ ベクトル（4-10） \\
\bkrouteicon{7} & はかる & ずれ（2-6）→ \R{標準偏差}{ひょう|じゅん|へん|さ}（3-14）→ \R{計測}{けい|そく}と\R{診断}{しん|だん}（9-4） \\
\end{bktable}
\bktabcaption{5\R{本}{ほん}の\R{路線}{ろ|せん}の、この\R{巻}{かん}の\R{大}{おお}きな\R{駅}{えき}}
\end{bkcenter}

\R{傾}{かたむ}きは「\R{変化}{へん|か}の\R{速}{はや}さ」、\R{指数関数}{し|すう|かん|すう}は「いまある\R{量}{りょう}に\R{比例}{ひ|れい}してふえる\R{量}{りょう}」、\R{座標}{ざ|ひょう}は「\R{場所}{ば|しょ}を2つの\R{数}{かず}で\R{表}{あらわ}す\R{方法}{ほう|ほう}」、\R{標準偏差}{ひょう|じゅん|へん|さ}は「ずれを1つの\R{数}{かず}で\R{表}{あらわ}したもの」です。\R{序章}{じょ|しょう}のばねの\R{実験}{じっ|けん}は、はかる\R{路線}{ろ|せん}（\bkrouteicon[0.9em]{7}）の\R{駅}{えき}です。はかる\R{路線}{ろ|せん}は、この\R{巻}{かん}でいちばんたくさんの\R{駅}{えき}に\R{止}{と}まります。たし\R{合}{あ}わせの\R{路線}{ろ|せん}（\bkrouteicon[0.9em]{2}）と\R{回転}{かい|てん}と\R{波}{なみ}の\R{路線}{ろ|せん}（\bkrouteicon[0.9em]{4}）は、この\R{巻}{かん}では\R{通}{とお}らず、\R{第}{だい}4\R{巻}{かん}（4-5・4-8）で\R{次}{つぎ}の\R{駅}{えき}に\R{着}{つ}きます。

\clearpage
%% ---------------- この巻でできるようになること ----------------
\bksection{goal}{この\R{巻}{かん}でできるようになること}

\R{第}{だい}3\R{巻}{かん}は、4つの\R{部}{ぶ}に\R{分}{わ}かれています。どの\R{部}{ぶ}も、ばねの\R{実験}{じっ|けん}で\R{見}{み}た5つの\R{段}{だん}（はかる・\R{表}{ひょう}・\R{点}{てん}・\R{線}{せん}・\R{式}{しき}）のどこかを、くわしくします。

\begin{bkfigure}
\begin{tikzpicture}[x=1mm, y=1mm,
  pt/.style={anchor=north west, draw=bkTheme, line width=1pt, fill=white, rounded corners=2.5mm,
    minimum width=64mm, minimum height=64mm, inner sep=0pt},
  pn/.style={anchor=north west, fill=bkTheme, text=bkOn, rounded corners=1.2mm, font=\sffamily\bfseries\small,
    inner xsep=2mm, inner ysep=1mm},
  pd/.style={anchor=north west, text width=57mm, align=left, font=\small, inner sep=0pt},
  px/.style={anchor=north, text=bkDeep, inner sep=0pt},
  pr/.style={anchor=north east, font=\sffamily\footnotesize, text=bkGray, inner sep=0pt}]
  \foreach \x/\y in {0/0, 68/0, 0/-70, 68/-70} {\node[pt] at (\x,\y) {};}
  % 第1部  % noruby
  \node[pn] at (3.5,-3.5) {\R{第}{だい}1\R{部}{ぶ}\enspace\R{関数}{かん|すう}とグラフ};
  \node[pr] at (61,-4.5) {3-1〜3-6};
  \node[px] at (32,-14) {$(3,\ 2)$\qquad$y=2x+5$};
  \node[pd] at (3.5,-23) {\bkthraitem{\R{場所}{ば|しょ}を2つの\R{数}{かず}（\R{座標}{ざ|ひょう}）で\R{表}{あらわ}せる。}\bkthraitem{\R{表}{ひょう}・グラフ・\R{式}{しき}を\R{行}{い}ったり\R{来}{き}たりできる。}\bkthraitem{グラフの\R{傾}{かたむ}きを\R{読}{よ}み、2\R{本}{ほん}の\R{直線}{ちょく|せん}の\R{交}{まじ}わる\R{点}{てん}を\R{求}{もと}められる。}};
  % 第2部  % noruby
  \node[pn] at (71.5,-3.5) {\R{第}{だい}2\R{部}{ぶ}\enspace\R{曲}{ま}がったグラフ};
  \node[pr] at (129,-4.5) {3-7〜3-10};
  \node[px] at (100,-13) {$y=\dfrac{a}{x}$\qquad$y=ax^2$};
  \node[pd] at (71.5,-23) {\bkthraitem{\R{反比例}{はん|ぴ|れい}と\R{放物線}{ほう|ぶつ|せん}のグラフを\R{読}{よ}める。}\bkthraitem{\R{式}{しき}のかっこをはずしたり、まとめたりできる。}\bkthraitem{ボールがその\R{高}{たか}さに\R{届}{とど}く\R{時刻}{じ|こく}を\R{計算}{けい|さん}できる。}};
  % 第3部  % noruby
  \node[pn] at (3.5,-73.5) {\R{第}{だい}3\R{部}{ぶ}\enspace\R{倍々}{ばい|ばい}とけた};
  \node[pr] at (61,-74.5) {3-11・3-12};
  \node[px] at (32,-83) {$2^{10}=1024$\qquad$\left(\dfrac{1}{2}\right)^3=\dfrac{1}{8}$};
  \node[pd] at (3.5,-95) {\bkthraitem{\R{倍々}{ばい|ばい}にふえる\R{量}{りょう}・\R{半分}{はん|ぶん}ずつへる\R{量}{りょう}を\R{計算}{けい|さん}できる。}\bkthraitem{けたのちがう\R{量}{りょう}を、けたの\R{数}{かず}（\R{対数}{たい|すう}）でくらべられる。}};
  % 第4部  % noruby
  \node[pn] at (71.5,-73.5) {\R{第}{だい}4\R{部}{ぶ}\enspace データと\R{統計}{とう|けい}};
  \node[pr] at (129,-74.5) {3-13〜3-17};
  \node[px] at (100,-84) {\R{平均}{へい|きん}\enspace±\enspace ばらつき};
  \node[pd] at (71.5,-95) {\bkthraitem{\R{平均}{へい|きん}と、ずれの\R{大}{おお}きさを1つの\R{数}{かず}で\R{言}{い}える。}\bkthraitem{ずれた\R{点}{てん}に、いちばん\R{合}{あ}う\R{直線}{ちょく|せん}を\R{引}{ひ}ける。}\bkthraitem{\R{偶然}{ぐう|ぜん}を\R{確率}{かく|りつ}で\R{考}{かんが}えられる。}};
\end{tikzpicture}
\bkfigcaption{4つの\R{部}{ぶ}と、できるようになること}
\end{bkfigure}

\begin{taisetsu}
どの\R{章}{しょう}でも、\R{見}{み}つけたきまりは「はかる → \R{表}{ひょう} → \R{点}{てん} → \R{線}{せん} → \R{式}{しき}」の\R{順}{じゅん}に\R{確}{たし}かめる。\par
\R{式}{しき}ができたら、\R{表}{ひょう}やグラフにもどって、\R{数}{かず}が\R{合}{あ}うかを\R{確}{たし}かめる。
\end{taisetsu}


\clearpage
\bksection{card}{この\R{巻}{かん}で\R{出会}{で|あ}う\R{新}{あたら}しい\R{記号}{き|ごう}}

\R{第}{だい}3\R{巻}{かん}では、\R{次}{つぎ}のような\R{新}{あたら}しい\R{記号}{き|ごう}に\R{出会}{で|あ}います。いまは\R{意味}{い|み}がわからなくても\R{大丈夫}{だい|じょう|ぶ}です。どの\R{記号}{き|ごう}も、\R{使}{つか}う\R{前}{まえ}に、\R{場面}{ば|めん}と\R{表}{ひょう}とグラフで、ていねいに\R{説明}{せつ|めい}します。\R{章}{しょう}の\R{中}{なか}の「\R{記号}{き|ごう}カード」は、\R{巻末}{かん|まつ}にまとめて\R{並}{なら}べてあります。

\begin{bkcenter}
\begin{bktable}{colspec={Q[c,m,wd=16mm] Q[l,m,wd=38mm] Q[l,m,wd=52mm] Q[c,m,wd=12mm]}, rowsep=1.5pt, cells={font=\small}, row{1}={font=\sffamily\bfseries\small}}
\R{記号}{き|ごう} & \R{読}{よ}み\R{方}{かた} & ひとことで\R{言}{い}うと & \R{章}{しょう} \\
$x_1$ & エックス・いち & 1\R{番}{ばん}めの $x$（かけ\R{算}{ざん}ではない） & 3-2 \\
$|x|$ & \R{絶対値}{ぜっ|たい|ち}エックス & 0からの\R{距離}{きょ|り} & 3-2 \\
$f(x)$ & エフ・エックス & \R{関数}{かん|すう} $f$ に $x$ を\R{入}{い}れた\R{値}{あたい} & 3-3 \\
$\alpha$、$\beta$ & アルファ、ベータ & 2\R{次方程式}{じ|ほう|てい|しき}の2つの\R{解}{かい} & 3-10 \\
$x^{\frac{1}{2}}$ & エックスの2\R{分}{ぶん}の1\R{乗}{じょう} & $\sqrt{x}$ のこと & 3-11 \\
$\log_ab$ & ログ & $a$ を\R{何乗}{なん|じょう}すると $b$ になるか & 3-12 \\
$\bar{x}$ & エックス・バー & \R{平均}{へい|きん} & 3-13 \\
$\sigma$ & シグマ & ばらつきの\R{大}{おお}きさ（\R{標準偏差}{ひょう|じゅん|へん|さ}） & 3-14 \\
$P(A)$ & ピー・エー & ことがら $A$ の\R{起}{お}こりやすさ（\R{確率}{かく|りつ}） & 3-17 \\
\end{bktable}
\bktabcaption{この\R{巻}{かん}で\R{出会}{で|あ}う\R{記号}{き|ごう}（くわしくはそれぞれの\R{章}{しょう}で）}
\end{bkcenter}

\R{記号}{き|ごう}は、\R{長}{なが}いことばを\R{短}{みじか}く\R{書}{か}くための\R{約束}{やく|そく}です。たとえば $f(3)$ は「\R{関数}{かん|すう} $f$ に3を\R{入}{い}れた\R{値}{あたい}」という\R{長}{なが}いことばを、4\R{文字}{も|じ}で\R{書}{か}いたものです。\R{約束}{やく|そく}を\R{知}{し}らないと\R{読}{よ}めませんが、\R{知}{し}ってしまえば、\R{世界}{せ|かい}じゅうの\R{人}{ひと}と\R{同}{おな}じ\R{意味}{い|み}でやりとりできます。

\R{同}{おな}じ\R{記号}{き|ごう}が、ちがう\R{意味}{い|み}で\R{使}{つか}われることもあります。$x_1$ はかけ\R{算}{ざん}の $x\times1$ ではありませんし、$f(x)$ は $f\times x$ ではありません。どちらも、その\R{章}{しょう}の「まちがえやすいところ」でくわしく\R{見}{み}ます。\R{第}{だい}5\R{巻}{かん}で\R{使}{つか}う\R{記号}{き|ごう}（$\lim$ や $e$ など）は、この\R{巻}{かん}では\R{使}{つか}いません。


\clearpage
%% ---------------- この巻の読み方 ----------------
\bksection{hon}{この\R{巻}{かん}の\R{読}{よ}み\R{方}{かた}}

\begin{bkfigure}
\begin{tikzpicture}[x=1mm, y=1mm]
  \bkthratalk{0}{sora}{ソラ}{\R{第}{だい}2\R{巻}{かん}では、\R{文字}{も|じ}の\R{式}{しき}や $F=kx$ を\R{作}{つく}ったね。\R{第}{だい}3\R{巻}{かん}では\R{何}{なに}をするの？}
  \bkthratalk{-21}{mina}{ミナ}{グラフがたくさん\R{出}{で}てくるみたい。$\log$ とか、\R{見}{み}たことのない\R{記号}{き|ごう}もある……。}
  \bkthratalk{-48}{delta}{デルタ}{だいじょうぶ。どの\R{章}{しょう}も、\R{場面}{ば|めん}を\R{見}{み}て、\R{表}{ひょう}を\R{作}{つく}って、グラフにして、\R{気}{き}づいてから\R{式}{しき}にするよ。\R{新}{あたら}しい\R{記号}{き|ごう}は、\R{使}{つか}う\R{前}{まえ}にかならず\R{説明}{せつ|めい}するよ。1\R{回}{かい}でわからなくても、\R{自分}{じ|ぶん}の\R{手}{て}で\R{点}{てん}をうってみれば、きっと\R{見}{み}えてくるよ。}
\end{tikzpicture}
\end{bkfigure}

\begin{bkfigure}
\begin{tikzpicture}[x=1mm, y=1mm]
  \bkthracard{(0,0)}{omoidasu}{bkGray}{①\enspace\R{思}{おも}い\R{出}{だ}そう}
    {\R{章}{しょう}のはじめに、\R{前}{まえ}に\R{学}{まな}んだことを1〜2\R{行}{ぎょう}で\R{思}{おも}い\R{出}{だ}します。\\[2pt]
     \bkthrapill{2-15}\enspace ばねの\R{力}{ちから}は $F=kx$\\[2pt]
     \R{札}{ふだ}の\R{章}{しょう}にもどれば、くわしく\R{書}{か}いてあります。}
  \bkthracard{(68,0)}{te}{bkNavy}{②\enspace\R{電卓}{でん|たく}・\R{表計算}{ひょう|けい|さん}を\R{使}{つか}おう}
    {\R{電卓}{でん|たく}、\R{表計算}{ひょう|けい|さん}（2-18）、プログラム（2-19）を\R{使}{つか}ってかまいません。\\
     \bkthraitem{おす\R{前}{まえ}に、\R{式}{しき}をノートに\R{書}{か}く。}\bkthraitem{\R{答}{こた}えのだいたいの\R{大}{おお}きさが\R{合}{あ}っているか、\R{確}{たし}かめる。}}
  \bkthracard{(0,-64)}{jibun}{bkOrange}{③\enspace ノートとグラフ}
    {\bkthraitem{\R{方眼}{ほう|がん}のノートを\R{使}{つか}う。\R{点}{てん}をうつ\R{問題}{もん|だい}が\R{多}{おお}い。}\bkthraitem{グラフの\R{軸}{じく}には、\R{量}{りょう}の\R{名前}{な|まえ}と\R{単位}{たん|い}を\R{書}{か}く。}\bkthraitem{\R{式}{しき}は1\R{行}{ぎょう}に1つ。「＝」をたてにそろえる。}}
  \bkthracard{(68,-64)}{reidai}{bkBlue}{④\enspace\R{例題}{れい|だい}から\R{章末問題}{しょう|まつ|もん|だい}へ}
    {\R{例題}{れい|だい}を\R{読}{よ}む → じぶんでとこう（\R{同}{おな}じ\R{形}{かたち}）→ \R{章末問題}{しょう|まつ|もん|だい}。\\[2pt]
     \R{章末問題}{しょう|まつ|もん|だい}の\R{形}{かたち}は、どれも\R{例題}{れい|だい}かじぶんでとこうで\R{一度}{いち|ど}\R{見}{み}せてあります。わからなかったら、\R{同}{おな}じ\R{形}{かたち}の\R{例題}{れい|だい}にもどろう。}
\end{tikzpicture}
\end{bkfigure}
